\textbf{1. Kiến thức cần nhớ}
\begin{itemize}
\item \textbf{Dữ liệu thống kê:} Gồm dữ liệu định tính (biểu thị bằng từ ngữ, chữ cái, kí hiệu) và dữ liệu định lượng (biểu thị bằng các con số, có thể đo đếm, tính toán).
\item \textbf{Bảng số liệu:} Sắp xếp dữ liệu theo hàng và cột để dễ tra cứu, so sánh. Thường gặp bảng thống kê, bảng tần số (số lần xuất hiện $n_i$ của mỗi giá trị), bảng tần suất (tỉ lệ phần trăm $f_i = \dfrac{n_i}{N} \cdot 100\%$ với $N$ là cỡ mẫu).
\item \textbf{Các loại biểu đồ thường gặp:}
\begin{itemize}
\item \textit{Biểu đồ cột (đơn, kép):} So sánh độ lớn của các đối tượng hoặc cùng một đối tượng ở các thời điểm khác nhau.
\item \textit{Biểu đồ đường (đoạn thẳng):} Biểu diễn xu thế thay đổi của dữ liệu theo thời gian.
\item \textit{Biểu đồ hình quạt tròn:} Biểu diễn tỉ lệ phần trăm (cơ cấu) của từng thành phần trong tổng thể (tổng bằng $100\%$).
\end{itemize}
\item \textbf{Số liệu bất thường (không hợp lí):} Là số liệu sai lệch rõ rệt so với quy luật chung hoặc mâu thuẫn với thực tế (như vượt quá giới hạn cho phép, tổng tỉ lệ khác $100\%$, sai đơn vị, tăng/giảm đột biến bất thường do lỗi ghi chép, nhập liệu).
\end{itemize}

\textbf{2. Phương pháp giải các dạng}

%% DANG 0D6?2-1
\textbf{Dạng 1. Đọc và phân tích thông tin trên bảng số liệu}
\begin{enumerate}
\item Đọc kĩ tiêu đề bảng, tên các hàng, tên các cột và đơn vị đo của từng đại lượng.
\item Đối chiếu giao điểm giữa hàng và cột để lấy đúng giá trị cần tìm.
\item Thực hiện các phép tính (tổng, hiệu, tỉ lệ phần trăm, so sánh lớn nhất/nhỏ nhất) và rút ra nhận xét theo yêu cầu đề bài.
\end{enumerate}
Lưu ý: Chú ý đơn vị tính ghi ở tiêu đề bảng (ví dụ: nghìn người, triệu đồng, tấn) để tránh sai lệch khi kết luận.

%% DANG 0D6?2-2
\textbf{Dạng 2. Đọc và phân tích thông tin trên biểu đồ}
\begin{enumerate}
\item Xác định loại biểu đồ, đọc tiêu đề, chú giải (màu sắc, kí hiệu) và đơn vị trên trục ngang, trục đứng.
\item Gióng từ đỉnh cột, điểm trên đường gấp khúc sang trục số (hoặc đọc trực tiếp số liệu/phần trăm ghi trên biểu đồ).
\item So sánh độ cao cột, độ dốc đường biểu diễn hoặc diện tích hình quạt để tính toán và nêu xu hướng tăng, giảm.
\end{enumerate}
Lưu ý: Với biểu đồ cột hoặc đường, cần kiểm tra gốc trục tung có bắt đầu từ $0$ hay không để tránh nhận xét sai lệch về mức độ chênh lệch.

%% DANG 0D6?2-3
\textbf{Dạng 3. Số liệu bất thường trên bảng số liệu}
\begin{enumerate}
\item Kiểm tra điều kiện thực tế và miền giá trị hợp lí của đại lượng (ví dụ: số người phải là số tự nhiên, điểm thi trong đoạn $[0; 10]$, không thể âm).
\item Kiểm tra tính nhất quán về tổng (tổng các phần bằng tổng chung, tổng tần suất bằng $100\%$) và so sánh các giá trị để tìm số liệu chênh lệch quá lớn.
\item Chỉ ra số liệu bất thường và giải thích nguyên nhân (nhập nhầm dấu phẩy, sai đơn vị, lỗi đếm).
\end{enumerate}
Lưu ý: Không nhầm lẫn giữa giá trị lớn nhất/nhỏ nhất hợp lí với số liệu bất thường; luôn đối chiếu với bối cảnh thực tế của bài toán.

%% DANG 0D6?2-4
\textbf{Dạng 4. Số liệu bất thường trên biểu đồ}
\begin{enumerate}
\item Quan sát tổng thể biểu đồ để phát hiện các điểm/cột nhô lên hoặc tụt xuống đột ngột, tách biệt hẳn so với xu hướng chung.
\item Kiểm tra tính hợp lí về mặt kĩ thuật: tỉ lệ hình vẽ có tương ứng với số liệu ghi trên biểu đồ không, tổng các phần trên biểu đồ quạt có bằng $100\%$ không.
\item Đối chiếu giá trị tại vị trí nghi ngờ với thực tế để kết luận số liệu bất thường và nêu lí do.
\end{enumerate}
Lưu ý: Trên biểu đồ hình quạt tròn, nếu phần trăm lớn hơn nhưng góc ở tâm lại vẽ nhỏ hơn thì biểu đồ biểu diễn sai tỉ lệ.
